\documentclass{article}
\usepackage[spanish]{babel}
\usepackage[utf8]{inputenc}
\usepackage{enumitem}
\usepackage{amsmath}
\usepackage{amssymb}
\usepackage{amsthm}
\usepackage{geometry}
\usepackage{setspace}
\usepackage{fancyhdr}
\geometry{letterpaper, margin=.7in}

% Edita rápido el ejercicio y la tarea
\newcommand{\ntarea}{1}

\begin{document}

\pagestyle{fancy}
\lhead{Probabilidad}
\rhead{Tarea \ntarea}

\large{Jesús Ezequiel Ramos Moreno.}
\begin{enumerate}
    \item[2.] Demuestre que si \(\{A_n\}\) es una colección de eventos en un mismo espacio de probabilidad todos con probabilidad 1, entonces \(\mathbb{P}(\bigcap_{n=1}^{\infty} A_n) = 1\).
    \begin{proof}
        Consideremos los complementos. Para cada \(n\), como \(P(A_n) = 1\), entnces \(A_n^c\)
    \end{proof}
\end{enumerate}

\end{document}